\chapter{Meniscus Tear}
\index{Meniscus Tear}

\section*{Definition and Overview}
A meniscus tear is a common cartilaginous injury of the knee joint, involving a rupture of one or both of the C-shaped fibrocartilaginous discs situated between the femoral condyles and the tibial plateau. The knee contains two menisci: the medial meniscus on the inner side and the lateral meniscus on the outer side. These structures are critical for the mechanical function of the knee, acting as shock absorbers, distributing axial loads, reducing friction during movement, and contributing to joint stability. Meniscus tears can occur as acute traumatic events, typically during sports involving twisting or pivoting motions, or as chronic degenerative changes in older individuals resulting from progressive wear and tear of the fibrocartilaginous matrix.

\section*{Epidemiology}
Meniscus tears represent one of the most frequently diagnosed knee injuries worldwide. Acute traumatic tears are highly prevalent in young, physically active populations (primarily aged 15 to 40), with a strong male predominance, often associated with pivoting and contact sports such as football, basketball, and netball. Degenerative tears are common in older adults, typically appearing in individuals over the age of 50. In these older cohorts, the prevalence of meniscus tears exceeds 30\%, often co-existing with osteoarthritis. The medial meniscus is injured significantly more often than the lateral meniscus, primarily because the medial meniscus is firmly anchored to the joint capsule and the medial collateral ligament (MCL), making it less mobile and more susceptible to shearing forces.

\section*{Aetiology and Pathophysiology}
The mechanism of a meniscus tear differs between traumatic and degenerative etiologies, influenced heavily by the vascular anatomy of the tissue.
\begin{itemize}
    \item \textbf{Traumatic mechanism:} Typically occurs when a twisting or rotational force is applied to a flexed, weight-bearing knee joint. This action shears the meniscus between the femur and the tibia. Traumatic tears are frequently associated with other knee ligament injuries, most notably anterior cruciate ligament (ACL) ruptures.
    \item \textbf{Degenerative mechanism:} Results from age-related degradation of the collagen network within the fibrocartilage. Over time, the meniscus loses its elastic properties and structural integrity, making it vulnerable to tearing from low-energy, routine activities such as squatting, kneeling, or rising from a chair.
    \item \textbf{Vascular zones:} The meniscus is divided into vascular zones which dictate its healing capacity. The outer third (the ``red-red'' or vascular zone) receives a rich blood supply from the perimeniscal capillary plexus, giving it the potential to heal, either spontaneously or with surgical repair. The inner two-thirds (the ``white-white'' or avascular zone) relies entirely on diffusion of synovial fluid for nutrition and has virtually no capacity for intrinsic repair.
\end{itemize}

\section*{Clinical Features}
The symptoms of a meniscus tear can vary depending on whether the tear is acute or degenerative, but typical clinical features include:
\begin{itemize}
    \item \textbf{Joint line pain:} Pain localised to the medial or lateral joint line, typically exacerbated by weight-bearing, twisting, or deep knee flexion.
    \item \textbf{Delayed swelling:} Development of a joint effusion (fluid on the knee) that typically appears 24 to 48 hours after an acute injury. This is in contrast to the rapid haemarthrosis (bleeding in the joint) seen in ACL tears.
    \item \textbf{Mechanical symptoms:} Intermittent clicking, popping, catching, or a sensation that the knee is giving way or slipping out of place during physical activity.
    \item \textbf{True joint locking:} An inability to fully extend or straighten the knee. This occurs classically with ``bucket-handle'' tears, where a large, displaced fragment of the meniscus flips into the intercondylar notch and physically blocks joint movement.
    \item \textbf{Localised tenderness:} Point tenderness upon palpation along the affected joint line during physical examination.
\end{itemize}

\section*{Differential Diagnosis}
Knee pain can arise from multiple structures, and a meniscus tear must be differentiated from:
\begin{itemize}
    \item \textbf{Knee osteoarthritis:} Presents with chronic joint pain, stiffness, and crepitus. Degenerative meniscus tears and osteoarthritis frequently co-exist, making it clinical diagnostic challenge to isolate the primary pain generator.
    \item \textbf{Collateral ligament sprain:} Injuries to the medial collateral ligament (MCL) or lateral collateral ligament (LCL), which present with localised joint line pain and swelling, often after a traumatic force.
    \item \textbf{Patellofemoral pain syndrome:} Pain located at the front of the knee, around or behind the kneecap, typically exacerbated by climbing stairs or prolonged sitting, without mechanical locking.
    \item \textbf{Osteochondritis dissecans:} An acquired lesion of the subchondral bone that can lead to cartilage delamination and loose bodies in the joint, mimicking mechanical symptoms of a tear.
    \item \textbf{Iliotibial band (ITB) friction syndrome:} Causes lateral knee pain, primarily in runners and cyclists, localized over the lateral femoral epicondyle rather than the joint line itself.
\end{itemize}

\section*{When to Seek Medical Care}
Individuals should consult a doctor if they experience persistent knee pain that does not improve with rest, recurrent joint swelling, or if they experience mechanical symptoms such as the knee locking, catching, or giving way.

\textbf{Contact your local emergency services for:}
\begin{itemize}
    \item An inability to bear any weight on the affected leg following an acute injury, or a visibly deformed knee joint (suggesting a fracture or joint dislocation).
    \item Signs of severe circulatory or nerve compromise in the leg, such as a cold, pale, or numb foot and calf.
    \item Severe, uncontrollable pain or massive, rapid swelling of the knee immediately following a high-impact trauma.
\end{itemize}

\section*{Diagnosis and Investigations}
A combination of a detailed clinical history, physical examination, and imaging is used to diagnose a meniscus tear.
\begin{itemize}
    \item \textbf{Physical examination tests:} Provocative clinical manoeuvres are performed to reproduce symptoms. These include \textbf{McMurray's test} (rotation of the tibia during knee extension to elicit a click or pain), the \textbf{Thessaly test} (internal and external rotation while standing on one slightly bent leg), and checking for direct joint line tenderness.
    \item \textbf{Magnetic Resonance Imaging (MRI):} The gold standard non-invasive imaging investigation. MRI is highly sensitive and specific, revealing the presence, location, size, and configuration of the tear (e.g., radial, longitudinal, horizontal, complex, or bucket-handle).
    \item \textbf{Plain radiographs (X-rays):} Recommended to rule out associated bony injuries, fractures, osteochondral lesions, or underlying osteoarthritis, although they cannot directly image the meniscus.
    \item \textbf{Diagnostic arthroscopy:} A minimally invasive surgical procedure that allows direct visualisation of the meniscus. While once common for diagnosis, it is now primarily reserved for cases where surgical intervention is concurrently planned.
\end{itemize}

\section*{Management}
Management is tailored to the patient's age, activity level, tear characteristics (location, stability), and the presence of mechanical symptoms.
\begin{itemize}
    \item \textbf{Conservative management:} The initial approach for small, stable tears, particularly degenerative tears without mechanical locking.
        \begin{itemize}
            \item \textbf{RICE protocol:} Rest, Ice application, Compression bandaging, and Elevation of the limb to control acute pain and swelling.
            \item \textbf{Pharmacotherapy:} Non-steroidal anti-inflammatory drugs (NSAIDs) such as ibuprofen or naproxen to reduce pain and inflammation.
            \item \textbf{Physical therapy:} A structured exercise programme focusing on quadriceps and hamstring strengthening, joint range of motion, and neuromuscular coordination.
        \end{itemize}
    \item \textbf{Surgical management:} Indicated for tears that fail conservative therapy, large or unstable tears, or tears presenting with acute joint locking.
        \begin{itemize}
            \item \textbf{Arthroscopic meniscal repair:} Suturing the torn meniscus back together. This is preferred in young patients with tears located in the vascular (red-red) zone to preserve the shock-absorbing function of the meniscus.
            \item \textbf{Arthroscopic partial meniscectomy:} Removing only the unstable, torn margins of the meniscus while preserving as much healthy tissue as possible. This is performed for tears located in the avascular (white-white) zone where healing is not possible.
        \end{itemize}
\end{itemize}

\section*{Complications}
A meniscus tear and its subsequent treatment can lead to long-term joint complications.
\begin{itemize}
    \item Accelerated knee osteoarthritis: Loss of meniscal tissue (especially following partial or total meniscectomy) reduces the contact area and increases stress on the articular cartilage, leading to premature joint degeneration.
    \item Persistent joint pain and recurrent effusions.
    \item Joint stiffness and permanent loss of full knee extension or flexion.
    \item \textbf{Surgical risks:} Post-operative joint infection (septic arthritis), deep vein thrombosis (DVT), or damage to local nerves (such as the saphenous nerve during medial meniscal repair).
\end{itemize}

\section*{Prognosis and Outcomes}
The prognosis for a meniscus tear varies based on the type of tear and the treatment pathway. Many patients with small degenerative tears achieve excellent pain relief and functional recovery with conservative management and physical therapy. For surgical cases, arthroscopic meniscal repair has a long-term success rate of approximately 80\% to 90\% in younger patients, preserving long-term joint health. Partial meniscectomy provides rapid short-term relief of pain and mechanical symptoms, but it carries a higher long-term risk of developing knee osteoarthritis, directly proportional to the amount of meniscal tissue removed.

\section*{Prevention and Self-Care}
While not all meniscus tears can be prevented, certain strategies can minimise the risk of injury.
\begin{itemize}
    \item Engage in regular strength training targeting the quadriceps, hamstrings, gluteals, and calf muscles to improve knee joint stability.
    \item Practice agility and balance drills to improve neuromuscular control during pivoting and landing manoeuvres.
    \item Perform a proper warm-up before participating in sports or vigorous physical activities.
    \item Maintain a healthy body weight to reduce the cumulative mechanical load on the knees.
    \item Wear appropriate, supportive footwear suitable for the specific sport or activity.
\end{itemize}

\section*{Sources and Further Reading}
The following reputable organisations provide further information and resources on meniscus tears:
\begin{itemize}
    \item Healthdirect Australia. \textit{Patient guide to meniscus tears, symptoms, and treatment options.} https://www.healthdirect.gov.au/
    \item OrthoInfo (American Academy of Orthopaedic Surgeons). \textit{Detailed clinical articles on meniscus tears and rehabilitation exercises.} https://orthoinfo.aaos.org/
    \item NHS (UK). \textit{Overview of meniscus tear symptoms, self-care, and surgical management.} https://www.nhs.uk/
    \item Australian Orthopaedic Association. \textit{Information on knee joint preservation and orthopaedic care.} https://aoa.org.au/
    \item MSD Manual (Professional Edition). \textit{Clinical reference guide for diagnosing and managing meniscal injuries.} https://www.msdmanuals.com/
\end{itemize}
